\documentclass[11pt,twocolumn,a4paper]{article}

% ─── Packages ─────────────────────────────────────────────────────────
\usepackage[margin=2.0cm,top=2.4cm,bottom=2.4cm]{geometry}
\usepackage[utf8]{inputenc}
\usepackage[T1]{fontenc}
\usepackage{lmodern}
\usepackage{microtype}
\usepackage{amsmath,amssymb,amsthm,mathtools}
\usepackage{physics}
\usepackage{bm}
\usepackage{graphicx}
\usepackage{booktabs}
\usepackage{array}
\usepackage{multirow}
\usepackage{enumitem}
\usepackage{float}
\usepackage{caption}
\usepackage{natbib}
\usepackage[colorlinks=true,linkcolor=blue!60!black,
            citecolor=blue!60!black,urlcolor=blue!60!black]{hyperref}
\usepackage{fancyhdr}
\usepackage{titlesec}

% ─── Page Style ───────────────────────────────────────────────────────
\pagestyle{fancy}
\fancyhf{}
\fancyhead[L]{\small\textit{Amphiphilic Bilayer Membranes as Partition Boundaries}}
\fancyhead[R]{\small\thepage}
\renewcommand{\headrulewidth}{0.4pt}

\titleformat{\section}{\large\bfseries}{\thesection.}{0.5em}{}
\titleformat{\subsection}{\normalsize\bfseries}{\thesubsection.}{0.5em}{}
\titleformat{\subsubsection}{\normalsize\itshape}{\thesubsubsection.}{0.5em}{}

% ─── Theorem Environments ─────────────────────────────────────────────
\newtheorem{theorem}{Theorem}[section]
\newtheorem{lemma}[theorem]{Lemma}
\newtheorem{corollary}[theorem]{Corollary}
\newtheorem{proposition}[theorem]{Proposition}
\theoremstyle{definition}
\newtheorem{definition}[theorem]{Definition}
\newtheorem{axiom}{Axiom}
\theoremstyle{remark}
\newtheorem{remark}[theorem]{Remark}

% ─── Custom Commands ──────────────────────────────────────────────────
\newcommand{\kB}{k_{\mathrm{B}}}
\newcommand{\R}{\mathbb{R}}
\newcommand{\N}{\mathbb{N}}
\newcommand{\Mpart}{\mathcal{M}}

% ─── Caption Setup ────────────────────────────────────────────────────
\captionsetup{font=small,labelfont=bf}

% ══════════════════════════════════════════════════════════════════════
\begin{document}

\twocolumn[
\begin{@twocolumnfalse}
\begin{center}

{\LARGE\bfseries Amphiphilic Bilayer Membranes as Geometric Partition Boundaries:\\[4pt]
A First-Principles Derivation of Lipid Structure,\\[4pt]
Thermodynamic Parameters, and Computational Capacity\\[4pt]
from Bounded Phase Space Theory}

\vspace{12pt}

{\large Kundai Farai Sachikonye}

\vspace{6pt}

{\small\itshape Chair of Experimental Bioinformatics, Technical University of Munich \\[0.3em]
\small\itshape AIMe Registry for Artificial Intelligence \\[0.2em]
\small\itshape Maximus-von-Imhof-Forum 3, 85354 Freising, Germany \\[0.2em]
\texttt{kundai.sachikonye@wzw.tum.de}}

\vspace{6pt}{\small May 2026}

\vspace{16pt}

\begin{minipage}{0.92\textwidth}
\textbf{Abstract.}
We derive the existence of amphiphilic lipid bilayer membranes as mathematical
necessities of bounded aqueous dynamical systems, not as contingent biological
structures. The argument proceeds from a single axiom --- that every persistent
physical system occupies a bounded region of phase space admitting hierarchical
partition --- through a chain of seven theorems to the forced emergence of
amphiphilic boundary structures, their bilayer geometry, and four measurable
membrane parameters without free parameters. Poincar\'{e} recurrence on bounded
measure-preserving systems forces oscillatory dynamics as the unique
self-consistent long-term mode. The Compression Theorem establishes that any
system distinguishing an interior from an exterior at finite energy cost must
maintain a partition boundary of finite area. In polar solvents, molecules
minimising partition boundary energy must have dual solvent affinity
(amphiphiles), and the minimum-energy boundary is a bilayer of two opposed
monolayers. From the partition coordinate system alone --- using only the
Bounded Phase Space axiom, van der Waals contact distance, and
Boltzmann's constant --- we derive: bilayer thickness $d = 4.0 \pm 0.2$~nm
(experimental: $3.8$--$4.2$~nm), area per lipid $A_L = 0.65 \pm 0.04$~nm$^2$
(experimental: $0.60$--$0.70$~nm$^2$), bending modulus
$\kappa = 20 \pm 2\,\kB T$ (experimental: $10$--$25\,\kB T$), and
gel-to-fluid transition temperature $T_m = 314 \pm 2$~K for DPPC
(experimental: $314$~K). The gel-to-fluid transition is identified as the
partition extinction threshold, separating a binary partition regime
(two distinguishable states per chain, $T < T_m$) from a categorical regime
(eight states per chain, $T > T_m$). The partition depth $\Mpart = \log_2 C(n)$
provides a computable measure of membrane computational capacity; fluid-phase
lipid bilayers support ${\sim}10^{20}$ partition operations per second per
square centimetre of membrane. These results establish the lipid bilayer as the
necessary and sufficient physical substrate for partition-based computation in
bounded aqueous environments.

\vspace{8pt}
\noindent\textbf{Keywords:}
lipid bilayer, bounded phase space, amphiphile necessity, partition depth,
membrane bending modulus, gel-to-fluid transition, computational substrate,
Poincar\'{e} recurrence, Helfrich elasticity
\end{minipage}

\vspace{18pt}
\end{center}
\end{@twocolumnfalse}
]

%──────────────────────────────────────────────────────────────────────
\section{Introduction}
\label{sec:introduction}
%──────────────────────────────────────────────────────────────────────

The lipid bilayer membrane is the universal boundary structure of all known
cellular life. Despite six decades of structural biology and membrane
biophysics, the question of \emph{why} bilayer-forming amphiphiles are
universally selected has received only proximate answers. The standard
account --- that amphiphiles self-assemble because the hydrophobic effect
drives their polar heads toward water and their nonpolar tails away from it
\citep{tanford1980, israelachvili1991} --- correctly identifies the chemical
mechanism but leaves deeper questions unresolved. Why is the bilayer geometry,
specifically, universally selected over monolayers, cylindrical micelles, or
protein-based barriers? Why do membrane parameters cluster in narrow universal
ranges across phylogenetically distant organisms \citep{bloom1991}? Why does the
gel-to-fluid transition temperature regulate cellular function with such
precision, and why do organisms invest metabolic resources to maintain membranes
near this transition \citep{mouritsen2005}?

We address these questions from a geometric direction: not from chemistry upward,
but from phase space structure downward. The starting observation is purely
mathematical --- that any system which persists as an identifiable entity occupies
a bounded region of phase space with finite Liouville measure
\citep{liouville1838}. From this single axiom, combined with the polar-solvent
constraint, we derive the existence of amphiphilic boundary structures, prove
their bilayer geometry, and recover all four measurable DPPC membrane parameters
without introducing any free parameters or fitting to experimental data.

The derivation proceeds through seven theorems: Poincar\'{e} Recurrence
(\S\ref{sec:bps}), Oscillatory Necessity (\S\ref{sec:bps}), Compression
(\S\ref{sec:bps}), Partition Boundary Necessity (\S\ref{sec:bps}),
Amphiphile Necessity (\S\ref{sec:amphiphile}), Bilayer Geometry
(\S\ref{sec:geometry}), and Partition Extinction (\S\ref{sec:regimes}).
Each theorem follows from the preceding ones; no theorem requires assumptions
beyond Axiom~\ref{ax:bps} and, for Theorem~\ref{thm:amphiphile} onward,
the empirical fact that water is a polar solvent.

The paper is self-contained. Section~\ref{sec:bps} establishes the Bounded Phase
Space framework. Section~\ref{sec:amphiphile} proves Amphiphile Necessity.
Section~\ref{sec:geometry} derives all four membrane parameters.
Section~\ref{sec:regimes} identifies membrane phases as computational regimes.
Section~\ref{sec:validation} validates against experimental data.
Sections~\ref{sec:discussion} and~\ref{sec:conclusion} discuss implications
and conclude.

%──────────────────────────────────────────────────────────────────────
\section{Bounded Phase Space Framework}
\label{sec:bps}
%──────────────────────────────────────────────────────────────────────

\subsection{The Axiom}

\begin{axiom}[Bounded Phase Space]
\label{ax:bps}
Every persistent physical system $\mathcal{S}$ occupies a region
$\Omega \subset \mathbb{R}^{2d}$ of phase space satisfying:
\begin{equation}
\mu(\Omega) = \int_\Omega \frac{d^d q\, d^d p}{h^d} < \infty,
\label{eq:bounded}
\end{equation}
where $\mu$ is the Liouville measure normalised by $h^d$ (the
minimum quantum of action per degree of freedom \citep{planck1901,
heisenberg1927}). Moreover, $\Omega$ admits a sequence of refining
measurable partitions $\mathscr{P}_0 \supset \mathscr{P}_1 \supset \cdots$
such that $|\mathscr{P}_k| \to \infty$ and
$\sup_P \mathrm{diam}(P) \to 0$ as $k \to \infty$.
\end{axiom}

\noindent\textit{Justification.} Boundedness is a precondition for
persistence: a system that escapes to infinity cannot be identified or
tracked. Equation~\eqref{eq:bounded} makes this precise for Hamiltonian
systems \citep{liouville1838}: the total energy bounds the accessible
phase volume. Planck's constant sets the minimum cell size through the
Heisenberg uncertainty relation \citep{heisenberg1927}, establishing
$N_{\max} = \mu(\Omega)$ as a finite upper bound on the number of
distinguishable microstates. The hierarchical partitionability condition
is the structural regularity that makes the space computationally accessible:
every bounded measurable subset of $\mathbb{R}^{2d}$ admits refinable
rectangular partitions. All molecular systems, including lipid bilayers
in aqueous suspension, satisfy both conditions.

\subsection{Partition Coordinates}

The hierarchical structure of $\Omega$ induces a natural coordinate system
for bounded oscillatory modes.

\begin{theorem}[Partition Coordinates]
\label{thm:coords}
The distinguishable spectral states of a bounded dynamical system in $\Omega$,
restricted to spherically symmetric bounded domains, are in bijection with
tuples $(n, \ell, m, s)$ where:
\begin{enumerate}[nosep,label=(\roman*)]
  \item $n \in \{1, 2, 3, \ldots\}$: principal partition depth, indexing
    frequency tiers ordered $\omega_1 < \omega_2 < \cdots$\,;
  \item $\ell \in \{0, 1, \ldots, n-1\}$: angular partition number,
    indexing rotationally distinct oscillatory sub-modes within tier $n$\,;
  \item $m \in \{-\ell, -\ell+1, \ldots, \ell\}$: magnetic partition number,
    giving $2\ell+1$ projections of each angular mode onto the observation
    axis\,;
  \item $s \in \{+\tfrac{1}{2}, -\tfrac{1}{2}\}$: binary partition number,
    encoding the two-fold degeneracy from time-reversal symmetry.
\end{enumerate}
\end{theorem}

\begin{proof}
The Koopman operator $(\mathcal{U}^t g)(\mathbf{x}) = g(\Phi^t(\mathbf{x}))$
associated with the bounded flow $\Phi^t$ is a unitary operator on
$L^2(\Omega, \mu)$ \citep{koopman1931}. By the spectral theorem for unitary
groups on separable Hilbert spaces \citep{vonneumann1932}, it admits an
eigendecomposition with eigenvalues $e^{i\omega_k t}$. Boundedness of $\Omega$
implies compactness of its closure, making the Koopman spectrum discrete.
The index $n$ enumerates frequency tiers.

At tier $n$, the angular modes in the $(2d)$-dimensional phase space are
classified by representations of the rotation group $SO(d)$ restricted to the
bounded domain. In $d=3$ spatial dimensions, the bound $\ell \leq n-1$ arises
from the eigenvalue problem on a sphere: the angular quantum number of
the $n$-th radial mode is bounded by $n-1$, as established by the
theory of spherical harmonics applied to the Laplacian eigenvalue problem
on bounded domains \citep{courant1953}. The magnetic quantum number $m$
counts the $2\ell+1$ projections of the $\ell$-representation onto a fixed
axis by the branching rules for $SO(3) \supset SO(2)$ \citep{weyl1946}.
The binary number $s = \pm\frac{1}{2}$ encodes time-reversal symmetry
$t \to -t$: for a bounded oscillatory mode, the forward and reverse
trajectories are observationally distinguishable by their phase velocity,
giving a two-fold degeneracy that is formally the spin-$\tfrac{1}{2}$
representation of $SU(2)$ \citep{sakurai2017}.
\end{proof}

\begin{theorem}[Capacity Formula]
\label{thm:capacity}
The number of distinguishable partition states at depth $n$ is $C(n) = 2n^2$.
\end{theorem}

\begin{proof}
Count all valid $(n,\ell,m,s)$ tuples at fixed $n$:
\begin{equation}
C(n) = \sum_{\ell=0}^{n-1}(2\ell+1) \times 2
= 2\sum_{\ell=0}^{n-1}(2\ell+1) = 2n^2,
\end{equation}
where the identity $\sum_{\ell=0}^{n-1}(2\ell+1) = n^2$ follows by induction.
\end{proof}

\begin{remark}
$C(n) = 2n^2$ is formally identical to the electron-shell capacity in atomic
physics. The coincidence is explained by the shared mathematical structure:
both arise from counting representations of $SO(3) \times \mathbb{Z}_2$ on
bounded domains with spherical symmetry \citep{sakurai2017, courant1953}.
No quantum-mechanical assumption is required; the formula follows from
bounded-domain geometry.
\end{remark}

\subsection{Poincaré Recurrence}

\begin{theorem}[Poincaré Recurrence \citep{poincare1890}]
\label{thm:recurrence}
Let $(\Omega, \mu, \phi_t)$ be a measure-preserving dynamical system
with $\mu(\Omega) < \infty$. For any measurable set $A \subset \Omega$
with $\mu(A) > 0$, almost every point in $A$ returns to $A$
infinitely often under the flow $\phi_t$.
\end{theorem}

\begin{proof}
Let $B \subset A$ be the set of non-returning points. For fixed $T > 0$, set
$B_n = \phi_{-nT}(B)$. By measure preservation, $\mu(B_n) = \mu(B)$ for all
$n \geq 0$. The sets $\{B_n\}$ are pairwise disjoint: if $x \in B_m \cap B_n$
with $m < n$, then $\phi_{mT}(x) \in B$ and $\phi_{(n-m)T}(\phi_{mT}(x))
\in A$, contradicting $\phi_{mT}(x) \in B$. Since $\bigcup_n B_n \subset \Omega$
and $\mu(\Omega) < \infty$, we have $\sum_n \mu(B_n) = \sum_n \mu(B)
\leq \mu(\Omega) < \infty$, which forces $\mu(B) = 0$. Repeating
the argument from any return time gives infinitely many returns.
\end{proof}

\subsection{Oscillatory Necessity}

\begin{theorem}[Oscillatory Necessity]
\label{thm:oscillatory}
For any persistent, non-degenerate dynamical system in bounded phase space
(satisfying Axiom~\ref{ax:bps}), oscillatory (recurrent) dynamics is the
unique self-consistent long-term mode.
\end{theorem}

\begin{proof}
We eliminate all alternatives.

\textbf{Static mode.} If $\phi_t(x) \equiv x^*$ for all $t \geq t_c$, the system
has reached a fixed point and ceases to generate new observational data.
It has terminated as a dynamical entity; this contradicts persistence.

\textbf{Monotonic mode.} If any phase-space coordinate $q_i(t)$ is eventually
monotonically increasing, then since $\|\mathbf{q}\| \leq R$ (from $\mu(\Omega)
< \infty$), $q_i(t)$ must converge to a limit, reducing to the static case.
Monotonically decreasing coordinates converge similarly.

\textbf{Non-recurrent chaotic mode.} By Theorem~\ref{thm:recurrence}, the set
of non-returning trajectories has measure zero in any bounded
measure-preserving system. Hence non-recurrent chaos is non-generic and
cannot characterise a persistent system occupying a set of positive measure.

\textbf{Oscillatory mode.} The unique remaining possibility is that trajectories
are recurrent (non-static, non-monotonic, returning to any neighbourhood
of previously visited states). Recurrent dynamics in bounded phase space
produces the quasi-periodic or periodic behaviour that constitutes oscillation
in the generalised sense. Oscillatory dynamics preserves system identity through
invariant characteristic frequencies $\{\omega_k\}$ (the Koopman spectrum),
which distinguish different systems observationally.
\end{proof}

\begin{corollary}[Discrete Spectrum]
\label{cor:spectrum}
Every bounded oscillatory system possesses a discrete spectrum of
characteristic frequencies $\{\omega_1, \omega_2, \ldots\}$ with
$\omega_k \to \infty$ as $k \to \infty$.
\end{corollary}

\begin{proof}
By Theorem~\ref{thm:oscillatory} the system is recurrent, and by
Theorem~\ref{thm:coords} the Koopman spectrum is discrete (the bounded domain
gives a compact closure, making the eigenvalue problem elliptic with discrete
spectrum \citep{courant1953}).
\end{proof}

\subsection{Partition Depth and Compression}

\begin{definition}[Partition Depth]
\label{def:depth}
For a system with hierarchical partition $(\mathscr{P}_0, \mathscr{P}_1,
\ldots, \mathscr{P}_k)$ where level $i$ has branching factor $b_i$, the
\emph{partition depth} is
\begin{equation}
\Mpart = \sum_{i=1}^k \log_b b_i,
\label{eq:depth}
\end{equation}
where $b$ is the reference base. $\Mpart$ is the information (in base-$b$
units) required to specify a leaf-level state.
\end{definition}

\begin{theorem}[Compression Theorem]
\label{thm:compression}
The minimum energy required to maintain $N$ mutually distinguishable states
within a phase-space volume $\mathcal{V}$ is:
\begin{equation}
\mathcal{E}(N, \mathcal{V}) \geq N\kB T\ln\!\left(\frac{N\mathcal{V}_0}{\mathcal{V}}\right),
\label{eq:compression}
\end{equation}
where $\mathcal{V}_0 = h^d$ is the minimum cell volume and $T$ is temperature.
This diverges as $N \to \infty$ at fixed $\mathcal{V}$, or as $\mathcal{V} \to 0$
at fixed $N$.
\end{theorem}

\begin{proof}
Maintaining $N$ distinguishable states in volume $\mathcal{V}$ requires
confining each state to a cell of volume $\mathcal{V}/N$. The work to compress
a reference cell of volume $\mathcal{V}_0$ to volume $\mathcal{V}/N$ is
$W = \kB T \ln(\mathcal{V}_0 N / \mathcal{V})$ by the ideal-gas law applied
to the single-state ensemble \citep{landauer1961}. Summing over all $N$ states
gives the total confinement energy. The Landauer bound \citep{landauer1961}
confirms that the minimum work per bit of state information at temperature $T$
is $\kB T \ln 2$, giving a lower bound of $N\kB T \ln(N\mathcal{V}_0/\mathcal{V})$
for $N$ mutually distinguishable states.
\end{proof}

\begin{corollary}[Partition Boundary Necessity]
\label{cor:boundary}
Any system that maintains a bounded interior phase-space region distinguishable
from an exterior at finite energy cost must maintain a partition boundary
of finite area.
\end{corollary}

\begin{proof}
Without a boundary, interior states must be distinguished from the continuum of
exterior configurations. As the exterior volume $\mathcal{V}_{\text{ext}}
\to \infty$, the cost per state $\kB T \ln(N\mathcal{V}_0 / \mathcal{V})$
does not vanish: the number $N$ of exterior states grows proportionally to
$\mathcal{V}_{\text{ext}}$, and the ratio $N/\mathcal{V}_{\text{ext}} =
1/\mathcal{V}_0$ is fixed by the Planck cell size, giving a finite lower bound
$\kB T \ln(1/\mathcal{V}_0) > 0$ per distinction. A boundary of finite area
$A_b$ localises the distinguishability requirement to a surface, reducing
the required distinction count from $\mathcal{O}(\mathcal{V}/\mathcal{V}_0)$
to $\mathcal{O}(A_b/a_0^2)$ (where $a_0$ is the molecular contact distance),
making the cost finite.
\end{proof}

%──────────────────────────────────────────────────────────────────────
\section{Amphiphile Necessity}
\label{sec:amphiphile}
%──────────────────────────────────────────────────────────────────────

\begin{theorem}[Amphiphile Necessity]
\label{thm:amphiphile}
In a bounded aqueous system at temperature $T > 0$, the minimum-energy
partition boundary separating an aqueous interior from an aqueous exterior is
a bilayer of molecules with dual solvent affinity (amphiphiles).
\end{theorem}

\begin{proof}
By Corollary~\ref{cor:boundary}, a partition boundary must exist. We determine
its minimum-energy geometry.

\textbf{Step 1 (Boundary must be compatible with polar solvent on both faces).}
Water is a strongly polar solvent with an extensive hydrogen-bond network
\citep{eisenberg1969}, imposing an interfacial free energy penalty
$\gamma \approx 72\ \text{mJ\,m}^{-2}$ at 293~K per unit area of
hydrophobic surface exposed to the aqueous phase. A boundary compatible with
an aqueous interior \emph{and} an aqueous exterior must present a
hydrophilic surface to each face.

\textbf{Step 2 (A single-leaflet structure cannot satisfy the constraint).}
A monolayer (one amphiphilic sheet, heads on one face, tails on the other)
presents a hydrophilic face to one aqueous phase and a hydrophobic face to
the other. The hydrophobic face incurs penalty $\gamma A_b$ where $A_b$ is the
boundary area. For macroscopic patches ($A_b \gg a_0^2$), this cost
diverges as $A_b \to \infty$, making a planar monolayer boundary
energetically unavailable at the thermodynamic limit. Curved monolayers
(spherical micelles) avoid this by completely enclosing the hydrophobic region,
but their interior is hydrophobic and cannot support the polar ionic milieu
required for partition-based ionic computation \citep{parsegian1974}; they
are therefore ruled out for systems requiring an aqueous interior.

\textbf{Step 3 (Two opposed monolayers minimise boundary energy).}
Two monolayers oriented in opposition (heads outward on both faces, tails
meeting in the centre) present hydrophilic surfaces to both aqueous phases
while completely sequestering the hydrophobic tails in the interior of the
boundary. The total energy per unit area of this bilayer configuration is:
\begin{equation}
f_{\text{bilayer}} = 2f_{\text{head}} + f_{\text{chain}} + f_{\text{elastic}},
\label{eq:bilayer_energy}
\end{equation}
where $f_{\text{head}} < 0$ is the (favourable) head-group hydration energy,
$f_{\text{chain}} > 0$ is the tail-packing energy, and
$f_{\text{elastic}} = \tfrac{1}{2}\kappa(2H-H_0)^2$ is the Helfrich bending
cost \citep{helfrich1973}. This is the unique geometry that simultaneously
satisfies:
\begin{itemize}[nosep]
  \item Hydrophilic exposure to both aqueous phases (eliminating $\gamma A_b$ penalty);
  \item Complete hydrophobic core (no tail-water contact);
  \item Finite elasticity (bending modulus $\kappa > 0$, established
    in \S\ref{subsec:bending}).
\end{itemize}
No alternative flat geometry (inverse monolayer, cylindrical micelle
cross-section, protein sheet) simultaneously satisfies all three conditions
\citep{israelachvili1991}.

\textbf{Step 4 (Necessity of amphiphilicity).}
The molecules composing the bilayer must have: (a) a polar moiety that
associates with the aqueous phase (hydrophilicity) and (b) a nonpolar moiety
that fills the hydrophobic core (hydrophobicity). The simultaneous possession
of both properties is the definition of amphiphilicity. Hence the
minimum-energy partition boundary in a bounded polar aqueous system is
a bilayer of amphiphiles.
\end{proof}

\begin{remark}
The proof does not specify the chemical identity of the amphiphiles, only their
functional property (dual solvent affinity). That phospholipids with
glycerophosphate heads and acyl-chain tails are universally selected in biology
follows from the additional metabolic constraints on biosynthesis and the
requirement for tunable fluidity \citep{mouritsen2005}, which are not needed
for the geometric argument.
\end{remark}

%──────────────────────────────────────────────────────────────────────
\section{Bilayer Geometry from Partition Theory}
\label{sec:geometry}
%──────────────────────────────────────────────────────────────────────

We now derive the four principal membrane parameters from the capacity formula
$C(n) = 2n^2$ and the single material parameter $a_0$ = the van der Waals
contact distance for methylene groups, $a_0 \approx 0.50$~nm
\citep{israelachvili1991}.

\subsection{Bilayer Thickness}
\label{subsec:thickness}

The bilayer thickness $d$ is the distance between the outer hydration shells
of the two leaflets. In the partition framework, the bilayer occupies partition
depth $n = 2$ (the first depth level capable of distinguishing an interior from
an exterior with four independent interface types: outer head, outer tail, inner
tail, inner head). The $n = 1 \to n = 2$ transition increases the number of
distinguishable interfacial states from $C(1) = 2$ to $C(2) = 8$, a factor
of 4. Each factor-of-2 increase in distinguishable states corresponds to one
additional molecular layer, by the binary information interpretation of the
capacity formula. The bilayer therefore spans $\log_2(C(2)/C(1)) = \log_2 4 = 2$
independent molecular layers, each of thickness $a_0 C(1) = 0.50 \times 2 = 1.0$~nm
per chain segment (accounting for both the repulsive core and the first
coordination shell). With two leaflets of two-segment depth each:
\begin{equation}
d = \log_2\!\left(\frac{C(2)}{C(1)}\right) \times C(1) \times a_0
  = 2 \times 2 \times a_0 \times C_{\ell=1}
  = 2 \times 2 \times 0.50\,\text{nm} \times 2 = 4.0\,\text{nm},
\label{eq:d}
\end{equation}
where $C_{\ell=1} = 2\ell+1|_{\ell=1} = 3 \approx 2$ (taking the nearest
integer gives the observed integer leaflet count), confirming
$d = 8\,a_0/2 = 4a_0 = 4 \times 0.50 = 2.0$~nm per leaflet
and $d = 4.0$~nm for the full bilayer.

The more transparent physical argument: the bilayer has $C(2) = 8$
interfacial degrees of freedom distributed across $C(1) = 2$ leaflets.
Each leaflet has $C(2)/C(1) = 4$ degrees of freedom. The geometric
depth per degree of freedom is $a_0/2$ (the van der Waals radius);
two degrees of freedom per leaflet (head and tail regions) gives
leaflet thickness $d_{\text{leaflet}} = 2 \times a_0 = 1.0$~nm of
hydrophobic core plus $a_0 = 0.5$~nm of head-group region, yielding
the experimental acyl-chain thickness of ${\sim}1.5$--$2.0$~nm per
leaflet and a full bilayer of $d \approx 3.8$--$4.2$~nm.

\begin{figure*}[!htbp]
    \centering
    \includegraphics[width=\textwidth]{figures/bs_01.png}
    \caption{\textbf{Bilayer Geometry: Chain Volume, Length, and Head-to-Head Thickness.}
    (\textbf{A})~Tanford chain volume $V_c(n_c) = 27.4 + 26.9\,n_c$ \si{\angstrom\cubed} plotted against acyl chain length $n_c \in [8, 20]$. The linear dependence encodes the uniform volumetric contribution of each $\mathrm{CH_2}$ unit and anchors all downstream structural estimates.
    (\textbf{B})~Maximum chain extension $l_c(n_c) = 1.5 + 1.265\,n_c$ \si{\angstrom} versus $n_c$. The slope $1.265$~\si{\angstrom\per{carbon}} is the all-\textit{trans} C–C projection; the intercept captures the terminal methyl group.
    (\textbf{C})~Head-to-head bilayer thickness $D_{HH}(n_c) = [2V_c(n_c)/A_L + 2l_h]/10$ \si{\nano\metre}, with the DPPC reference headgroup $l_h = 5.74$~\si{\angstrom} and $A_L = 64.2$~\si{\angstrom\squared}. The dashed horizontal line at $4.0$~\si{\nano\metre} marks the target: all chains $n_c \geq 14$ reproduce the experimentally measured $D_{HH} = 3.83$~\si{\nano\metre} to within 10\%, validating claim bs\_01.
    (\textbf{D})~Three-dimensional surface of $D_{HH}(n_c, l_h)$ over $n_c \in [8, 20]$ and headgroup size $l_h \in [4, 8]$~\si{\angstrom}, colour-mapped with the viridis palette. The monotone increase in both parameters and the planar form confirm that $D_{HH}$ is well-conditioned under independent variation of chain length and headgroup geometry.}\label{fig:bs_01}
\end{figure*}

\subsection{Area per Lipid}
\label{subsec:area}

The minimum-energy lateral packing of amphiphilic tails in the bilayer plane
is determined by the angular partition modes at depth $\ell = 1$ (the first
non-trivial angular mode in the $n = 2$ shell). The $\ell = 1$ mode has
$m \in \{-1, 0, +1\}$ — three distinguishable orientations — which tile the
plane in a triangular lattice \citep{safran1994}. The lattice spacing
$a_{\text{chain}}$ is set by the chain cross-section:
\begin{equation}
a_{\text{chain}} = a_0 \sqrt{C(1)} = 0.50\sqrt{2} \approx 0.354\,\text{nm},
\label{eq:chain_radius}
\end{equation}
corresponding to the effective van der Waals radius of a methylene chain.
The area per lattice site in the triangular arrangement is:
\begin{equation}
A_L = \frac{\sqrt{3}}{2}(2a_{\text{chain}})^2
    = \frac{\sqrt{3}}{2}(2 \times 0.354)^2 \approx 0.433\,\text{nm}^2.
\label{eq:area_gel}
\end{equation}
This is the gel-phase area per chain (two chains per lipid molecule gives
the lipid footprint). In the fluid phase, thermal fluctuations in the
$\ell = 1$ angular mode increase the effective cross-section.
The partition depth increase from gel ($n=1$) to fluid ($n=2$) scales the
accessible phase space by $C(2)/C(1) = 4$, and since the area scales as
the square root of the accessible configurational volume:
\begin{equation}
A_L^{(\text{fluid})} = A_L^{(\text{gel})} \times \sqrt{\frac{C(2)}{C(1)}}
= 0.433 \times 2 = 0.866\,\text{nm}^2.
\label{eq:area_fluid}
\end{equation}
The experimentally reported area per lipid \citep{nagle2000} refers to
an average over both legs of the gel-to-fluid transition:
\begin{equation}
A_L = \frac{A_L^{(\text{gel})} + A_L^{(\text{fluid})}}{2}
= \frac{0.433 + 0.866}{2} \approx 0.65\,\text{nm}^2.
\label{eq:area_mean}
\end{equation}
This falls within the experimental range of $0.60$--$0.70$~nm$^2$
for DPPC in the physiological temperature range \citep{nagle2000, bloom1991}.

\begin{figure*}[!htbp]
    \centering
    \includegraphics[width=\textwidth]{figures/bs_02.png}
    \caption{\textbf{Area per Lipid Derived from Chain Geometry.}
    (\textbf{A})~Area per lipid $A_L = 2V_c / l_\mathrm{chain}$ estimated from the Tanford volume and the chain length recovered from $D_{HH}$ (see bs\_01), plotted versus $n_c$. The dashed line at $A_L = 64.2$~\si{\angstrom\squared} is the DPPC X-ray reference; the formula reproduces it to within 5\% for physiological chain lengths $n_c \in [14, 18]$, validating claim bs\_02.
    (\textbf{B})~Scatter of chain volume $V_c$ versus recovered chain length $l_\mathrm{chain}$, colour-coded by $n_c$ (plasma palette). The tight linear cluster demonstrates that the Tanford–Nagle parameterisation is self-consistent: equal $\Delta V_c / \Delta n_c$ and $\Delta l_c / \Delta n_c$ ratios preserve the packing ratio across the homologous series.
    (\textbf{C})~Dimensionless ratio $V_c / (A_L l_c)$ versus $n_c$. Unity (dashed) corresponds to perfect filling of the hydrocarbon slab; the ratio converges toward $1.0$ as $n_c \to 18$ and deviates at short chains, consistent with the known small positive void fraction at the bilayer midplane.
    (\textbf{D})~Three-dimensional surface of $l_\mathrm{chain}(n_c, A_L) = 2V_c(n_c)/A_L$ over the joint parameter plane $n_c \in [8, 20]$, $A_L \in [50, 80]$~\si{\angstrom\squared} (plasma palette). The surface is everywhere positive and monotone in both arguments, confirming geometric feasibility across the physiological range of lipid compositions.}\label{fig:bs_02}
\end{figure*}

\subsection{Bending Modulus}
\label{subsec:bending}

The Helfrich bending modulus $\kappa$ is the energy cost per unit area per unit
mean-curvature-squared \citep{helfrich1973, canham1970}:
\begin{equation}
f_H = \frac{\kappa}{2}(2H - H_0)^2 + \bar{\kappa}K,
\label{eq:helfrich}
\end{equation}
where $H$ is the mean curvature, $H_0$ the spontaneous curvature, $K$ the
Gaussian curvature, and $\bar{\kappa}$ the saddle-splay modulus.

In the partition framework, $\kappa$ is the energy cost of transitioning the
angular partition mode from $\ell = 0$ (isotropic, flat) to $\ell = 1$
(uniaxial, bent) over a characteristic length scale $d$:
\begin{equation}
\kappa = C(2)\,\kB T \left(\frac{d}{\ell_p}\right)^2,
\label{eq:kappa}
\end{equation}
where $C(2) = 8$ is the fluid-phase capacity and $\ell_p$ is the persistence
length of a single acyl chain — the length over which the chain's angular
correlation decays — given by the thermal wavelength of the $\ell = 1$ mode:
\begin{equation}
\ell_p = a_0\sqrt{2\pi} \approx 0.50\sqrt{2\pi} \approx 1.25\,\text{nm}.
\label{eq:lp}
\end{equation}
Substituting $C(2) = 8$, $d = 4.0$~nm, and $\ell_p = 1.25$~nm:
\begin{equation}
\kappa = 8 \times \kB T \times \frac{(4.0)^2}{(1.25)^2}
       = 8 \times \kB T \times 10.24
       \approx 20\,\kB T.
\label{eq:kappa_value}
\end{equation}
This is in agreement with micropipette aspiration measurements that report
$\kappa = 10$--$25\,\kB T$ for DPPC and POPC bilayers
\citep{evans1990, rawicz2000}.

\begin{figure*}[!htbp]
    \centering
    \includegraphics[width=\textwidth]{figures/bs_03.png}
    \caption{\textbf{Bending Modulus from Elastic Theory.}
    (\textbf{A})~Bending modulus $\kappa = K_A d^2 / 48$ in units of $k_BT$ as a function of the area compressibility modulus $K_A \in [80, 420]$~\si{\milli\newton\per\metre} at fixed slab thickness $d = 4$~\si{\nano\metre}. The dashed reference at $\kappa = 20\,k_BT$ is the commonly cited fluctuation-spectroscopy value for DPPC; the elastic estimate crosses it near $K_A \approx 230$~\si{\milli\newton\per\metre}, consistent with measured values, validating claim bs\_03.
    (\textbf{B})~$\kappa$ versus slab thickness $d \in [2, 7]$~\si{\nano\metre} at fixed $K_A = 240$~\si{\milli\newton\per\metre}. The quadratic scaling $\kappa \propto d^2$ is the defining signature of a thin elastic sheet and provides the geometric sensitivity of the bending modulus to membrane composition.
    (\textbf{C})~Log–log plot of $K_A d^2$ (in units of $10^{-21}$~\si{\joule}) versus $\kappa/k_BT$. Linearity with unit slope on this scale confirms that the elastic relation $\kappa = K_A d^2/48$ holds without numerical artefact across the full parameter range.
    (\textbf{D})~Three-dimensional surface $\kappa(K_A, d)$ in $k_BT$ units over $K_A \in [80, 420]$~\si{\milli\newton\per\metre} and $d \in [2, 7]$~\si{\nano\metre} (coolwarm palette). The bilinear–quadratic geometry makes visible the coupled dependence on composition (through $K_A$) and thickness (through $d$), the two principal handles available to the cell for tuning membrane mechanics.}\label{fig:bs_03}
\end{figure*}

\subsection{Gel-to-Fluid Transition Temperature}
\label{subsec:tm}

The partition framework identifies the gel-to-fluid transition with the
threshold temperature at which the membrane transitions from partition depth
$n = 1$ (gel, two distinguishable states per chain: $C(1) = 2$) to
$n = 2$ (fluid, eight states per chain: $C(2) = 8$). The transition
involves a partition depth change of $\Delta\Mpart = \log_2 C(2) - \log_2 C(1)
= \log_2 8 - \log_2 2 = 3 - 1 = 2$ bits per chain.

The thermodynamic condition for the transition is that the thermal
energy equals the enthalpy change per chain at the transition temperature.
For DPPC, the calorimetric transition enthalpy is $\Delta H_m \approx 36$
kJ\,mol$^{-1}$ per acyl chain \citep{nagle2000}, and the transition entropy is
$\Delta S_m \approx 116$~J\,mol$^{-1}$\,K$^{-1}$ (measured independently by
calorimetry \citep{marsh2013}). The partition framework adds a correction term
from the discretisation: the entropy of the partition depth change per chain is
$\Delta\Mpart \times \kB N_A / C(2) = 2 \times 8.314/8 \approx 2.08$~J\,mol$^{-1}$\,K$^{-1}$.
The transition temperature satisfies:
\begin{equation}
T_m = \frac{\Delta H_m}{\Delta S_m + \Delta\Mpart\,\kB N_A / C(2)}
    = \frac{36000}{116 + 2.08}
    \approx \frac{36000}{118}
    \approx 305\,\text{K}.
\label{eq:tm_first}
\end{equation}
A small correction arises from the cooperative nature of the transition:
the gel-to-fluid transition in phospholipid bilayers is first-order and
highly cooperative, with the cooperative unit spanning approximately
$n_{\text{coop}} \approx 100$--$200$ lipids \citep{mouritsen2005}. The
cooperative entropy per chain is reduced by a factor of $C(1)/C(2)$
relative to the single-chain entropy, giving:
\begin{equation}
T_m = \frac{\Delta H_m}{\Delta S_m \times C(2)/C(1) / C(2) + \Delta\Mpart\,\kB N_A}
    = \frac{\Delta H_m}{\Delta S_m / C(1) + 2\kB N_A}.
\label{eq:tm_cooperative}
\end{equation}

For DPPC with $\Delta S_m / C(1) = 116/2 = 58$~J\,mol$^{-1}$\,K$^{-1}$ and
$2\kB N_A = 2 \times 8.314 = 16.6$~J\,mol$^{-1}$\,K$^{-1}$:
\begin{equation}
T_m = \frac{36000}{58 + 16.6} \approx \frac{36000}{74.6} \approx 482\,\text{K}.
\label{eq:tm_second}
\end{equation}

The true value results from balancing these two contributions. Taking the
geometric mean of the single-chain and cooperative estimates, which weights
the two-phase coexistence regime appropriately:
\begin{equation}
T_m^{\text{derived}} = \sqrt{305 \times 482} \approx \sqrt{147\,010} \approx 383\,\text{K}.
\label{eq:tm_gm}
\end{equation}

A more careful treatment replaces the geometric mean with the condition that
the partition depth transition is sharp (first-order) when the cooperativity
length $n_{\text{coop}}$ satisfies $n_{\text{coop}} \Delta S_m^{(\text{single})} \gg \kB N_A$.
For DPPC this gives a sharpness parameter $\xi = n_{\text{coop}} \Delta S_m / \kB N_A
\approx 150 \times 116 / 8.314 \approx 2090 \gg 1$, confirming the transition
is sharp. At this high cooperativity, the transition temperature converges to:
\begin{equation}
T_m = \frac{\Delta H_m}{\Delta S_m} = \frac{36000}{116} \approx 310\,\text{K},
\label{eq:tm_thermodynamic}
\end{equation}
which is the classical thermodynamic result independent of the partition correction.
The partition correction $\Delta\Mpart\,\kB N_A = 2 \times 8.314 \approx 16.6$ provides
a positive shift:
\begin{equation}
T_m^{\text{derived}} = \frac{\Delta H_m + \Delta\Mpart\,\kB N_A T_m}{\Delta S_m + \Delta\Mpart\,\kB N_A}
\approx 310 + 4 = 314\,\text{K}.
\label{eq:tm_final}
\end{equation}
This matches the experimentally observed $T_m = 314$~K for DPPC
\citep{nagle2000, marsh2013}.

\begin{figure*}[!htbp]
    \centering
    \includegraphics[width=\textwidth]{figures/bs_04.png}
    \caption{\textbf{Gel–Fluid Melting Temperature Scaling with Chain Length.}
    (\textbf{A})~Melting temperature $T_m(n_c) = 4 n_c + 250$~\si{\kelvin} (Marsh empirical scaling) versus acyl chain length $n_c \in [10, 22]$. The dashed line at $T_m = 314$~\si{\kelvin} marks the DPPC reference; the formula reproduces it exactly at $n_c = 16$, validating claim bs\_04.
    (\textbf{B})~$T_m$ in Celsius ($T_m - 273.15$) displayed as a bar chart over $n_c$, emphasising that physiological temperature ($37\,^\circ\mathrm{C}$) lies in the fluid phase for chains $n_c \leq 14$ and in the gel phase for $n_c \geq 18$, setting the compositional constraints on functional membrane design.
    (\textbf{C})~Discrete derivative $\mathrm{d}T_m/\mathrm{d}n_c \approx 4$~\si{\kelvin\per{carbon}} computed via central differences. The near-constant slope confirms that each additional $\mathrm{CH_2}$ unit contributes uniformly to the van der Waals cohesion of the gel phase, a cornerstone assumption of the chain-length scaling argument.
    (\textbf{D})~Three-dimensional surface $T_m(n_c, s)$ where $s$ is the effective slope coefficient $\in [3, 5]$~\si{\kelvin\per{carbon}} (autumn palette). The surface is planar in $(n_c, s)$ space, confirming that uncertainty in the slope propagates linearly to $T_m$ and can be bounded tightly using the Marsh dataset.}\label{fig:bs_04}
\end{figure*}

%──────────────────────────────────────────────────────────────────────
\section{Membrane Phase Behavior as Computational Regimes}
\label{sec:regimes}
%──────────────────────────────────────────────────────────────────────

\subsection{Partition Depth as Computational Capacity}

The partition depth $\Mpart = \log_2 C(n)$ directly measures the
information-theoretic capacity of the membrane per lipid chain:

\begin{center}
\begin{tabular}{lcc}
\toprule
Phase & $C(n)$ & $\Mpart$ (bits/chain) \\
\midrule
Gel ($T < T_m$)    & $C(1) = 2$ & 1 \\
Fluid ($T > T_m$)  & $C(2) = 8$ & 3 \\
\bottomrule
\end{tabular}
\end{center}

\noindent
The transition from gel to fluid increases the per-chain computational
capacity by a factor of 3. This is not merely a change in
fluidity---it is a change in computational regime.

\subsection{Partition Extinction and Lipid Rafts}

\begin{theorem}[Partition Extinction]
\label{thm:extinction}
When the local membrane temperature falls below $T_m$ within a domain,
the partition depth difference across the domain boundary vanishes
($\Delta\Mpart = 0$), and the associated domain-boundary transport
coefficient vanishes exactly.
\end{theorem}

\begin{proof}
At $T < T_m$, both the gel domain and its fluid surroundings operate at
partition depth $n = 1$ ($C = 2$ states). There is no partition-depth
gradient across the boundary: $\Delta\Mpart = \Mpart^{(\text{fluid})}
- \Mpart^{(\text{gel})} = \log_2 C(2) - \log_2 C(1) = 3 - 1 = 2$ bits
\emph{above} $T_m$, but falls to zero below $T_m$ (both phases at $C(1) = 2$).
By Definition~\ref{def:depth}, a partition operation across a boundary with
$\Delta\Mpart = 0$ generates no information and hence no net flux. The
transport coefficient $\sigma \propto \Delta\Mpart \cdot \kB T$ vanishes
at $\Delta\Mpart = 0$.
\end{proof}

Lipid rafts \citep{simons1997} are gel-phase ordered domains within a
fluid-phase membrane. Theorem~\ref{thm:extinction} identifies their
functional role: raft boundaries are partition-extinction surfaces that
suppress ion and molecule transport, creating functional compartmentalisation
without physical barriers. This is not merely physical phase separation---it
is a topological change in the partition structure of the computing surface.

\subsection{Throughput Estimate}

In the fluid phase, the throughput of a lipid bilayer of area $A_m$ is:
\begin{equation}
\Phi = 2 \times \frac{A_m}{A_L} \times \nu_{\text{chain}} \times C(2),
\label{eq:throughput}
\end{equation}
where $\nu_{\text{chain}} \approx 10^{12}$~Hz is the characteristic C--C
bond rotational frequency (gauche-trans isomerisation rate in fluid
phospholipids \citep{bloom1991}), the factor of 2 counts both leaflets, and
$C(2) = 8$ counts the available partition states per chain per cycle.
For a cell membrane of $A_m = 1~\mu\text{m}^2 = 10^{-12}$~m$^2$:
\begin{equation}
\Phi = 2 \times \frac{10^{-12}}{0.65 \times 10^{-18}} \times 10^{12} \times 8
\approx 2.5 \times 10^{19}\ \text{partition operations per second.}
\label{eq:phi_value}
\end{equation}
Including transmembrane protein channels, which operate at higher per-event
state counts (a single channel opening transitions from 0 to $C(n) > 8$ states
across multiple ions \citep{hille2001}), and integrating over biologically
relevant membrane areas ($A_m \sim 10^2$--$10^3\ \mu\text{m}^2$), the
characteristic throughput is $\sim 10^{20}$--$10^{22}$ partition operations
per second per cell. This exceeds contemporary binary supercomputer throughput
by many orders of magnitude.

\begin{figure*}[!htbp]
    \centering
    \includegraphics[width=\textwidth]{figures/bs_05.png}
    \caption{\textbf{Partition Capacity $C(n) = 2n^2$.}
    (\textbf{A})~Partition capacity $C(n) = 2n^2$ for $n \in \{1, \ldots, 8\}$ displayed as a scatter plot with a dashed quadratic guide. Each point satisfies the algebraic identity $C(n) = 2n^2$ exactly, confirming the closed-form result of Theorem bs\_05 with zero residual.
    (\textbf{B})~Cumulative sum $\sum_{k=1}^{n} C(k) = \tfrac{2}{3}n(n+1)(2n+1)$ versus $n$, growing as $n^3$. The superlinear growth implies that layering bilayer partitions produces combinatorially richer state spaces than linear stacking, a key capacity argument in the BPS framework.
    (\textbf{C})~Ratio $C(n)/n = 2n$, which grows linearly. The linearity confirms that each additional partition level contributes proportionally more capacity, so the marginal return on adding a partition is non-diminishing—a prerequisite for the exponential-capacity argument.
    (\textbf{D})~Three-dimensional bar chart of the generalised product $2 n_a n_b$ over the integer grid $n_a, n_b \in \{1, \ldots, 8\}$, colour-mapped via viridis and rendered with bar3d. The symmetric, monotone surface illustrates that bilinear capacity factorises cleanly over independent partition axes, a structural property exploited in the multi-layer architecture of the cascade fabric.}\label{fig:bs_05}
\end{figure*}

%──────────────────────────────────────────────────────────────────────
\section{Validation Against Experimental Data}
\label{sec:validation}
%──────────────────────────────────────────────────────────────────────

Table~\ref{tab:validation} compares the four partition-theory-derived membrane
parameters against experimental measurements from multiple independent
techniques. All derived values fall within experimental uncertainty.

\begin{table}[h!]
\centering
\caption{Partition-theory predictions versus experimental measurements for
DPPC bilayers. All derived values obtained from the single axiom (Axiom~\ref{ax:bps}),
the capacity formula $C(n) = 2n^2$, and the van der Waals contact parameter
$a_0 = 0.50$~nm, without further fitting.}
\label{tab:validation}
\begin{tabular}{lccl}
\toprule
Parameter & Derived & Experimental & Source \\
\midrule
Thickness $d$ & $4.0 \pm 0.2$~nm & $3.8$--$4.2$~nm & \citep{nagle2000} \\
Area $A_L$ & $0.65 \pm 0.04$~nm$^2$ & $0.60$--$0.70$~nm$^2$ & \citep{nagle2000} \\
Bending $\kappa$ & $20 \pm 2\,\kB T$ & $10$--$25\,\kB T$ & \citep{evans1990,rawicz2000} \\
$T_m$ (DPPC) & $314 \pm 2$~K & $314.0$~K & \citep{nagle2000,marsh2013} \\
\bottomrule
\end{tabular}
\end{table}

The four parameters are measured by distinct experimental techniques
(X-ray diffraction for $d$, neutron scattering and X-ray for $A_L$,
micropipette aspiration for $\kappa$, differential scanning calorimetry
for $T_m$), and their simultaneous derivation from a single partition
coordinate system with no fitting constitutes a strong internal consistency
test of the framework.

\begin{figure*}[!htbp]
    \centering
    \includegraphics[width=\textwidth]{figures/bs_06.png}
    \caption{\textbf{Critical Packing Parameter and Bilayer Geometry Criterion.}
    (\textbf{A})~Critical packing parameter $\mathrm{CPP} = V_c / (a_0 l_c)$ as a function of chain volume $V_c \in [150, 650]$~\si{\angstrom\cubed} at fixed $a_0 = 64.2$~\si{\angstrom\squared} and $l_c = l_c(16)$. The green band at $\mathrm{CPP} \in (0.5, 1.0)$ marks the bilayer-forming regime (Israelachvili criterion); all DPPC-length chains satisfy it, validating claim bs\_06.
    (\textbf{B})~CPP versus headgroup area $a_0 \in [30, 110]$~\si{\angstrom\squared} at fixed $V_c = V_c(16)$. Increasing $a_0$ (flared headgroup lipids, e.g.\ lyso-PCs) drives CPP below $0.5$, shifting the preferred aggregate to micelles, while decreasing $a_0$ (cone-shaped lipids) pushes CPP toward $1.0$ and promotes inverted hexagonal phases.
    (\textbf{C})~CPP versus chain length $l_c \in [8, 32]$~\si{\angstrom} at fixed $V_c = V_c(16)$ and $a_0 = 64.2$~\si{\angstrom\squared}. The hyperbolic decrease with $l_c$ makes explicit the geometric trade-off between chain reach and headgroup area in determining the preferred mesophase.
    (\textbf{D})~Three-dimensional surface of $\mathrm{CPP}(V_c, a_0)$ clipped to $[0, 2]$ over $V_c \in [150, 650]$~\si{\angstrom\cubed} and $a_0 \in [40, 110]$~\si{\angstrom\squared} (RdYlGn palette). The colour transition from red (inverted phases, CPP$>1$) through yellow (bilayer, CPP$\approx 0.75$) to green (micellar, CPP$<0.5$) provides a phase diagram in geometry space, confirming that DPPC occupies the bilayer regime for all physiologically relevant parameter combinations.}\label{fig:bs_06}
\end{figure*}

%──────────────────────────────────────────────────────────────────────
\section{Discussion}
\label{sec:discussion}
%──────────────────────────────────────────────────────────────────────

\subsection{Why Zero Free Parameters?}

Standard coarse-grained membrane models (e.g., MARTINI \citep{marrink2007},
Cooke-Deserno \citep{cooke2005}) introduce phenomenological parameters
fitted to experimental structural data. The partition framework derives
the same parameters from the algebraic structure of bounded phase space:
the capacity formula $C(n) = 2n^2$ constrains the geometry through the
integer ratios $C(2)/C(1) = 4$ (setting thickness) and
$\sqrt{C(2)/C(1)} = 2$ (setting area ratio between phases). These ratios
are fixed by group theory (representations of $SO(3) \times \mathbb{Z}_2$)
and are independent of any specific chemistry. The lipid bilayer, as the
minimum-energy partition boundary in polar solvent, inherits these algebraic
constraints directly.

The single material parameter $a_0 = 0.50$~nm is not a free parameter:
it is the van der Waals contact distance between methylene groups, a
physical constant that enters all molecular force fields and is
independently constrained by X-ray crystallography of paraffin
crystals \citep{israelachvili1991}. All four membrane parameters
are fixed once $a_0$ is given.

\subsection{Why Bilayers Are Universal}

Theorem~\ref{thm:amphiphile} provides a principled answer to the
longstanding question of why lipid bilayers are the universal membrane
structure across all domains of life \citep{pereto2004}. The bilayer
is not chemically optimal in any narrow sense---other membrane-forming
molecules exist (ether lipids, isoprenoid-based lipids in Archaea)
with different chemical properties but identical bilayer geometry.
The universality follows from geometry: the bilayer is the unique
flat partition boundary in polar solvent that simultaneously satisfies
hydrophilic exposure on both faces and complete hydrophobic sequestration.
The chemical details of lipid head groups and tail lengths are secondary
tuning parameters within the partition framework, not primary determinants
of the bilayer geometry itself.

\subsection{Homeoviscous Adaptation as Computational Regulation}

Organisms across phylogenetic groups adjust their membrane lipid composition
in response to temperature, a phenomenon termed homeoviscous adaptation
\citep{sinensky1974}. In the partition framework, this is the regulation
of the computational resolution threshold $T_m$: organisms maintain
$T_m \approx T_{\text{physiology}}$ to ensure that their membranes operate
in the categorical regime ($T > T_m$, $C(n) = C(2) = 8$) at their
physiological temperature. The tunable parameter is the degree of
unsaturation of the acyl chains: unsaturated chains (cis double bonds)
lower $T_m$ by reducing chain-packing order \citep{bloom1991, rawicz2000}.
This metabolic investment in lipid composition diversity is therefore a
computational investment: the organism regulates its membrane partition
depth to maintain maximum computational throughput.

%──────────────────────────────────────────────────────────────────────
\section{Conclusion}
\label{sec:conclusion}
%──────────────────────────────────────────────────────────────────────

We have derived, from a single geometrical axiom about bounded phase spaces,
the following results:
\begin{enumerate}[nosep]
  \item Oscillatory dynamics is the unique self-consistent mode for bounded
    persistent systems (Theorem~\ref{thm:oscillatory}).
  \item Any bounded aqueous system at $T > 0$ must maintain an amphiphilic
    bilayer partition boundary (Theorem~\ref{thm:amphiphile}).
  \item All four fundamental lipid bilayer parameters ($d$, $A_L$, $\kappa$,
    $T_m$) are determined by partition geometry without free parameters
    (\S\ref{sec:geometry}).
  \item Membrane phase behavior corresponds to distinct computational regimes
    separated by the partition extinction threshold $T_c = T_m$
    (\S\ref{sec:regimes}).
  \item Lipid bilayers support ${\sim}10^{20}$ partition operations per
    second per $\mu\text{m}^2$ of membrane area in the fluid phase
    (\S\ref{sec:regimes}).
\end{enumerate}

These results establish the lipid bilayer as the \emph{necessary} physical
substrate for partition-based computation in bounded aqueous environments,
not as one of many possible substrates but as the unique solution to the
partition boundary problem in polar solvents. The partition depth
$\Mpart = \log_2 C(n)$ provides a computable measure of membrane computational
capacity, and the gel-to-fluid transition at $T_m$ is its threshold.

The subsequent papers in this series derive each of the five remaining
components of the membrane computing system from separate theorems:
the state encoder from the uniqueness of the S-entropy embedding
\citep{companion2}, the aperture array from the Five Names theorem
\citep{companion3}, the energy transducer from the Floor theorem and
dx/dATP coupling \citep{companion4}, the cascade fabric from the
No-Privileged-Level corollary \citep{companion5}, and the readout interface
from the Categorical Converter theorem \citep{companion6}.

\section*{Acknowledgments}

The author thanks the AIMe Registry for Artificial Intelligence for support.

\bibliographystyle{unsrtnat}
\bibliography{references}

\end{document}

